\documentclass[10pt,a4paper]{article}
\usepackage{fontspec}
\setmainfont{NotoSans-Regular.ttf}[BoldFont=NotoSans-Bold.ttf,ItalicFont=NotoSans-Italic.ttf,BoldItalicFont=NotoSans-BoldItalic.ttf]
\setsansfont{NotoSans-Regular.ttf}[BoldFont=NotoSans-Bold.ttf]
\setmonofont{NotoSansMono-Regular.ttf}[Scale=0.87]
\usepackage[russian]{babel}
\usepackage[margin=20mm,headheight=15pt]{geometry}
\usepackage{xcolor}
\definecolor{ink}{HTML}{202138}
\definecolor{accent}{HTML}{7140BD}
\definecolor{muted}{HTML}{606575}
\definecolor{panel}{HTML}{F5F2FA}
\usepackage{microtype}
\usepackage{parskip}
\usepackage{titlesec}
\usepackage{fancyhdr}
\usepackage{fvextra}
\usepackage[most]{tcolorbox}
\usepackage{xurl}
\usepackage[colorlinks=true,urlcolor=accent,linkcolor=accent,unicode,pdftitle={ResearcherOS StarterKIT},pdfauthor={ResearcherOS}]{hyperref}
\usepackage{enumitem}
\setlist{nosep}
\newtcolorbox{promptbox}{colback=panel,colframe=accent!25,boxrule=0.4pt,arc=3pt,left=9pt,right=9pt,top=7pt,bottom=7pt,before skip=8pt,after skip=10pt}
\fvset{fontsize=\small,breaklines=true,breakanywhere=true,breaksymbolleft={},breaksymbolright={},baselinestretch=1.1}
\titleformat{\section}{\Large\bfseries\color{ink}}{}{0pt}{}
\titleformat{\subsection}{\normalsize\bfseries\color{accent}}{}{0pt}{}
\titlespacing*{\section}{0pt}{0pt}{14pt}
\titlespacing*{\subsection}{0pt}{12pt}{5pt}
\setcounter{secnumdepth}{0}
\pagestyle{fancy}
\fancyhf{}
\fancyhead[L]{\small\color{muted}ResearcherOS / StarterKIT}
\fancyhead[R]{\small\color{muted}Практическое руководство}
\fancyfoot[L]{\scriptsize\color{muted}Числовой wCTR · регрессия · CPU}
\fancyfoot[R]{\small\thepage}
\renewcommand{\headrulewidth}{0.3pt}
\setlength{\emergencystretch}{3em}
\color{ink}
\begin{document}
\begin{titlepage}
\vspace*{12mm}
{\color{accent}\large RESEARCHEROS}\par
\vspace{14mm}
{\fontsize{38}{44}\selectfont\bfseries StarterKIT}\par
\vspace{7mm}
{\Large Первое исследование\\с AI-агентом}\par
\vspace{12mm}
{\large Прогноз взвешенной доли кликов\\на данных рекламы велосипедов}\par
\vspace{12mm}
\begin{promptbox}
Восемь шагов: подключение проекта, дерево, карточка, постановка,
запуск, результаты, редактура и два режима Arbor.
\end{promptbox}
\vspace{7mm}
Читайте действия, копируйте готовое сообщение в указанный чат и проверяйте результат в конце каждого шага. Терминальные команды выполняются отдельно.
\vfill
\href{https://github.com/ZoyaV/researcheros_starter_kit}{github.com/ZoyaV/researcheros\_starter\_kit}\par
\small Набор данных: Sell me this! / 9\,923 строки / 2023-2024\par
Версия руководства: 9 ноября 2026 г.
\end{titlepage}
\tableofcontents
\vspace{10mm}
\subsection*{До начала}
Нужны Git, Git LFS, Python 3.10 или новее, ResearcherOS и настроенный CLI-агент. Сам эксперимент работает на CPU. Подписка или API выбранного агента нужны отдельно.

Ветка исходного набора - main. Дерево вы создадите сами; исследовательские материалы хранит ResearcherOS отдельно от кода.
\subsection*{Правила сравнения}
Предсказываем числовой wCTR, не вероятность клика. Основная метрика - RMSE, меньше лучше. Постоянный прогноз сравнивается с деревом, затем с улучшенными вариантами.

Фиксированное разделение по датам: train - 6\,093 строки, validation - 1\,895, test - 1\,935. Подбор выполняется на validation. Итоговый test открывается один раз в конце шага 8. Учебный порог улучшения - 1\%; он не заменяет проверку статистической значимости.

\href{https://github.com/ZoyaV/researcheros_starter_kit/blob/main/data/README.md}{Полное описание данных и протокола в репозитории.}
% WALKTHROUGH_CONTENT
\end{document}
